% Źródła: materialy/2026-11-09/IMG_7102--IMG_7115.jpeg.
\poczatekwykladu{6}

\begin{fakt}[Niezależność zmiennych dyskretnych]
\label{fakt:niezaleznosc-zmiennych-dyskretnych}
Niech $X_1,\ldots,X_N$ będą zmiennymi losowymi o rozkładach dyskretnych.
Są one niezależne wtedy i tylko wtedy, gdy dla każdego wyboru punktów
$x_i\in S_{X_i}$, $i=1,\ldots,N$, zachodzi równość
\[
  \PP(X_1=x_1,\ldots,X_N=x_N)=\prod_{i=1}^{N}\PP(X_i=x_i).
\]
\end{fakt}
\noindent\emph{Dowód pozostawiono jako ćwiczenie.}

\section{Rozkłady absolutnie ciągłe}
\label{sec:rozklady-absolutnie-ciagle}

\begin{definicja}
\label{def:rozklad-absolutnie-ciagly}
Rozkład $\mu$ na $\R$ nazywamy \emph{absolutnie ciągłym}, jeżeli istnieje
funkcja borelowska $f\colon\R\to[0,\infty)$ taka, że dla każdego
$A\in\B(\R)$
\[
  \mu(A)=\int_A f(x)\,dx
        =\int_{\R}\mathbf{1}_A(x)f(x)\,dx.
\]
Funkcję $f$ nazywamy \emph{gęstością} rozkładu $\mu$.
Jeżeli $\mu_X$ jest absolutnie ciągły, mówimy, że zmienna losowa $X$
ma rozkład absolutnie ciągły. Jego gęstość oznaczamy przez $f_X$.
\end{definicja}

\begin{fakt}
Jeżeli $X$ ma rozkład absolutnie ciągły, to $S_X=\varnothing$.
\end{fakt}

\begin{proof}
Dla dowolnego $x\in\R$ mamy
\[
  \PP(X=x)=\int_{\{x\}} f_X(u)\,du=0,
\]
ponieważ singleton $\{x\}$ ma miarę Lebesgue'a równą zeru.
\end{proof}

\begin{uwaga}[Gęstość a rozkład]
\label{uw:gestosc-a-rozklad}
\begin{enumerate}
  \item Gęstość jest wyznaczona z dokładnością do zbioru miary
    Lebesgue'a zero. Dokładniej, dla nieujemnej funkcji borelowskiej $f$
    zachodzi równoważność
    \[
      f=f_X\ \text{prawie wszędzie}
      \quad\Longleftrightarrow\quad
      \mu_X(A)=\int_A f(x)\,dx
      \quad\text{dla każdego }A\in\B(\R).
    \]
  \item Każda funkcja borelowska $f\colon\R\to[0,\infty)$, dla której
    $\int_{\R}f(x)\,dx=1$, wyznacza rozkład absolutnie ciągły przez
    wzór $\mu(A)=\int_A f(x)\,dx$ dla $A\in\B(\R)$.
\end{enumerate}
\end{uwaga}

\subsection{Przykłady rozkładów absolutnie ciągłych}

\begin{przyklad}[Rozkład jednostajny]
\label{przykl:rozklad-jednostajny}
Dla $a<b$ rozkład jednostajny $U([a,b])$ ma gęstość
\[
  f_X(x)=\frac{1}{b-a}\,\mathbf{1}_{[a,b]}(x).
\]
\end{przyklad}

\begin{przyklad}[Rozkład wykładniczy]
\label{przykl:rozklad-wykladniczy}
Dla $\lambda>0$ rozkład wykładniczy $\operatorname{Exp}(\lambda)$
ma gęstość
\[
  f_X(x)=\lambda e^{-\lambda x}\mathbf{1}_{(0,\infty)}(x).
\]
Istotnie, całka z tej funkcji wynosi
\[
  \int_{\R} f_X(x)\,dx
  =\int_0^\infty\lambda e^{-\lambda x}\,dx
  =\lim_{b\to\infty}\bigl(1-e^{-\lambda b}\bigr)=1.
\]
\end{przyklad}

\begin{figure}[H]
  \centering
  \includegraphics[width=12cm]{gestosci-jednostajna-wykladnicza.pdf}
  \caption{Gęstość rozkładu jednostajnego z przykładu~\ref{przykl:rozklad-jednostajny}
    (z lewej) i wykładniczego z przykładu~\ref{przykl:rozklad-wykladniczy}
    (z prawej).}
  \label{fig:gestosci-jednostajna-wykladnicza}
\end{figure}

\begin{przyklad}[Rozkład normalny]
\label{przykl:rozklad-normalny}
Rozkład Gaussa $N(m,\sigma^2)$, gdzie $m\in\R$ i $\sigma>0$,
ma gęstość
\[
  f_X(x)=\frac{1}{\sigma\sqrt{2\pi}}
  \exp\!\left(-\frac{(x-m)^2}{2\sigma^2}\right),
  \qquad x\in\R.
\]
Standardowy rozkład normalny $N(0,1)$ ma gęstość
\[
  f_X(x)=\frac{1}{\sqrt{2\pi}}e^{-x^2/2},
  \qquad x\in\R.
\]
\end{przyklad}

\begin{figure}[H]
  \centering
  \includegraphics[width=12cm]{gestosci-normalne.pdf}
  \caption{Gęstości normalne z przykładu~\ref{przykl:rozklad-normalny}.
    Po lewej $N(0,1)$, po prawej dwie gęstości o wspólnym $m$.
    Węższa krzywa odpowiada mniejszemu $\sigma$.}
  \label{fig:gestosci-normalne}
\end{figure}

\begin{przyklad}[Rozkład gamma]
\label{przykl:rozklad-gamma}
Dla $\alpha>0$ i $\beta>0$ rozkład gamma ma gęstość
\[
  f_X(x)=
  \begin{cases}
    \dfrac{\beta^\alpha}{\Gamma(\alpha)}\,x^{\alpha-1}e^{-\beta x},
      & x>0,\\[5pt]
    0, & x\leq 0.
  \end{cases}
\]
Funkcja gamma jest dana wzorem
\[
  \Gamma(\alpha)=\int_0^\infty x^{\alpha-1}e^{-x}\,dx,
  \qquad \alpha>0.
\]
Dla $n=0,1,2,\ldots$ mamy $\Gamma(n+1)=n!$.
Przy $\alpha=1$ otrzymujemy gęstość rozkładu wykładniczego
$\operatorname{Exp}(\beta)$.
\end{przyklad}

\begin{przyklad}[Czas oczekiwania na tramwaj]
\label{przykl:oczekiwanie-na-tramwaj}
Niech $X$ oznacza czas oczekiwania na tramwaj w minutach i niech
$X\sim\operatorname{Exp}(\lambda)$, gdzie $\lambda=1/15$.
\begin{enumerate}[label=\alph*)]
  \item Prawdopodobieństwo oczekiwania dokładnie $10$ minut wynosi
    $\PP(X=10)=0$, ponieważ rozkład jest absolutnie ciągły.
  \item Prawdopodobieństwo oczekiwania dłużej niż $10$ minut wynosi
    \[
      \PP(X>10)=\int_{10}^{\infty}\lambda e^{-\lambda x}\,dx
      =e^{-10\lambda}=e^{-2/3}.
    \]
  \item Prawdopodobieństwo oczekiwania od $10$ do $20$ minut wynosi
    \[
      \PP(10\leq X\leq 20)=\int_{10}^{20}\lambda e^{-\lambda x}\,dx
      =e^{-10\lambda}-e^{-20\lambda}=e^{-2/3}-e^{-4/3}.
    \]
\end{enumerate}
\end{przyklad}

\section{Dystrybuanta}
\label{sec:dystrybuanta}

\begin{definicja}
\label{def:dystrybuanta}
\emph{Dystrybuantą} rozkładu $\mu_X$ rzeczywistej zmiennej losowej $X$
nazywamy funkcję $F_X\colon\R\to[0,1]$ daną wzorem
\[
  F_X(t)=\PP(X\leq t)=\mu_X((-\infty,t]),\qquad t\in\R.
\]
\end{definicja}

\begin{przyklad}[Dystrybuanta rozkładu wykładniczego]
\label{przykl:dystrybuanta-wykladnicza}
Niech $X\sim\operatorname{Exp}(\lambda)$, gdzie $\lambda>0$.
Dla $t\leq 0$ mamy $F_X(t)=0$. Dla $t>0$ otrzymujemy
\[
  F_X(t)=\int_0^t\lambda e^{-\lambda x}\,dx
        =\bigl[-e^{-\lambda x}\bigr]_{0}^{t}
        =1-e^{-\lambda t}.
\]
\end{przyklad}

\begin{przyklad}[Dystrybuanta rozkładu dwumianowego]
\label{przykl:dystrybuanta-dwumianowa}
Niech $X\sim B(2,\tfrac12)$. Z wzoru z przykładu~\ref{przykl:rozklad-bernoulliego}
otrzymujemy $\PP(X=0)=\tfrac14$, $\PP(X=1)=\tfrac12$
oraz $\PP(X=2)=\tfrac14$. Zatem
\[
  F_X(t)=
  \begin{cases}
    0, & t<0,\\
    \tfrac14, & 0\leq t<1,\\
    \tfrac34, & 1\leq t<2,\\
    1, & t\geq 2.
  \end{cases}
\]
\end{przyklad}

\begin{figure}[H]
  \centering
  \includegraphics[width=12.5cm]{dystrybuanty-wykladnicza-dwumianowa.pdf}
  \caption{Dystrybuanty z przykładów~\ref{przykl:dystrybuanta-wykladnicza}
    i~\ref{przykl:dystrybuanta-dwumianowa}.
    Na wykresie rozkładu wykładniczego przyjęto $\lambda=1$.}
  \label{fig:dystrybuanty-wykladnicza-dwumianowa}
\end{figure}

\begin{twierdzenie}[Własności dystrybuanty]
\label{tw:wlasnosci-dystrybuanty}
Niech $F_X$ będzie dystrybuantą zmiennej losowej $X$.
\begin{enumerate}[label=(\roman*)]
  \item Funkcja $F_X$ jest niemalejąca.
  \item Zachodzą równości
    \[
      \lim_{t\to-\infty}F_X(t)=0,\qquad
      \lim_{t\to\infty}F_X(t)=1.
    \]
  \item Funkcja $F_X$ jest prawostronnie ciągła.
  \item Dla każdego $t_0\in\R$ istnieje granica lewostronna
    \[
      F_X(t_0^-)=\lim_{t\uparrow t_0}F_X(t)=\PP(X<t_0).
    \]
  \item Funkcja $F_X$ jest nieciągła w $t_0$ wtedy i tylko wtedy, gdy
    $\PP(X=t_0)>0$. Ponadto
    \[
      \PP(X=t_0)=F_X(t_0)-F_X(t_0^-).
    \]
\end{enumerate}
\end{twierdzenie}

\begin{proof}
\begin{enumerate}[label=(\roman*)]
  \item Dla $t_1\leq t_2$ mamy $(-\infty,t_1]\subseteq(-\infty,t_2]$,
    więc z monotoniczności miary wynika $F_X(t_1)\leq F_X(t_2)$.
  \item Z monotoniczności $F_X$ wystarczy rozpatrzyć wartości w punktach
    $-n$ oraz $n$. Korzystając z ciągłości miary $\mu_X$, otrzymujemy
    \begin{align*}
      \lim_{n\to\infty}F_X(-n)
        &=\mu_X\!\left(\bigcap_{n=1}^\infty(-\infty,-n]\right)
          =\mu_X(\varnothing)=0,\\
      \lim_{n\to\infty}F_X(n)
        &=\mu_X\!\left(\bigcup_{n=1}^\infty(-\infty,n]\right)
          =\mu_X(\R)=1.
    \end{align*}
  \item Jeżeli $t_n\downarrow t_0$, to przedziały $(-\infty,t_n]$
    maleją do $(-\infty,t_0]$. Zatem
    \[
      \lim_{n\to\infty} F_X(t_n)
        =\mu_X\!\left(\bigcap_{n=1}^{\infty}(-\infty,t_n]\right)
        =F_X(t_0).
    \]
  \item Jeżeli $t_n\uparrow t_0$ i $t_n<t_0$, to przedziały
    $(-\infty,t_n]$ rosną do $(-\infty,t_0)$. Stąd
    \[
      F_X(t_0^-)=\lim_{n\to\infty}F_X(t_n)
        =\mu_X((-\infty,t_0))=\PP(X<t_0).
    \]
  \item Odejmując prawdopodobieństwa, dostajemy
    \[
      \PP(X=t_0)=\PP(X\leq t_0)-\PP(X<t_0)=F_X(t_0)-F_X(t_0^-).
    \]
    Z prawostronnej ciągłości wynika, że nieciągłość zachodzi dokładnie
    wtedy, gdy ta różnica jest dodatnia. \qedhere
\end{enumerate}
\end{proof}

\begin{definicja}
\label{def:rozklad-ciagly}
Rozkład $\mu_X$ nazywamy \emph{ciągłym}, gdy jego dystrybuanta $F_X$
jest ciągła.
\end{definicja}

Przykład ciągłej dystrybuanty z odcinkami stałości przedstawia
rysunek~\ref{fig:dystrybuanta-cantora}.

\begin{figure}[H]
  \centering
  \includegraphics[width=8.5cm]{dystrybuanta-cantora.pdf}
  \caption{Piąte przybliżenie dystrybuanty Cantora. W granicy otrzymujemy
    funkcję ciągłą, stałą na każdym z przedziałów usuwanych przy budowie
    zbioru Cantora.}
  \label{fig:dystrybuanta-cantora}
\end{figure}

\begin{uwaga}
Każda funkcja niemalejąca, prawostronnie ciągła i o granicach $0$
w $-\infty$ oraz $1$ w $+\infty$ jest dystrybuantą pewnego rozkładu.
Są to warunki \textup{(i)--(iii)} twierdzenia~\ref{tw:wlasnosci-dystrybuanty}.
\end{uwaga}
